\section{Introduction}
Large language models can still fail on deterministic micro-behaviors (e.g., basic comparison or parity) even when they remain useful on broader tasks.
In this setting, post-hoc repair methods should satisfy three practical requirements:
(i) repair the targeted behavior reliably, (ii) limit collateral drift on unrelated prompts, and (iii) provide auditable evidence for what was actually fixed.

We study this problem under a frozen-model assumption and focus on inference-time interventions.
The resulting protocol, CertiPatch, combines constrained optimization with explicit artifact verification, so that reported outcomes are replayable from saved run artifacts.
The citable release artifact for this manuscript is archived at \href{https://doi.org/10.5281/zenodo.18541322}{10.5281/zenodo.18541322}.
This places the method in the overlap between parameter-efficient adaptation methods such as LoRA and prompt tuning~\cite{hu2022lora,lester2021prompt},
localized model-editing ideas~\cite{meng2022rome,meng2023memit},
and test-driven repair loops.

\paragraph{Contributions.}
\begin{itemize}[leftmargin=*,nosep]
\item \textbf{Train-certify-verify protocol.} We provide an end-to-end repair pipeline that emits certificate artifacts and enforces fail-closed verification.
\item \textbf{Scope-aware certification.} We report exact certification on enumerable domains and coverage-bounded certification on a large stratified domain.
\item \textbf{Compositionality evidence.} We evaluate six composition conditions and quantify interference/order effects.
\item \textbf{Compute-aware claim discipline.} We restrict comparative claims to completed cells and keep the submission single-seed by design.
\end{itemize}
